% Generated by book/tools/notebook_to_tex.py. Do not edit by hand.

\chapter[TF-IDF]{TF-IDF로 만나는 첫 벡터}

\markboth{TF-IDF}{TF-IDF}

\chaptermeta{01_tfidf/01_tfidf.ipynb}{https://colab.research.google.com/github/yoon-gu/neuqes-101/blob/master/01_tfidf/01_tfidf.ipynb}{텍스트를 숫자 벡터로 바꾸는 첫 관문}



\index{1TF-IDF@TF-IDF}

\index{1CountVectorizer@CountVectorizer}

\index{1TfidfVectorizer@TfidfVectorizer}

\index{1sparse matrix@sparse matrix}

\index{1vocabulary@vocabulary}

\index{0텍스트 벡터화@텍스트 벡터화}

\index{0단어 빈도@단어 빈도}

\index{0희귀도 가중치@희귀도 가중치}

\index{0어휘@어휘}

\index{0희소 행렬@희소 행렬}


\index{1Bag of Words@Bag of Words}

\index{1BoW@BoW}

\index{1n-gram@n-gram}

\index{1token\_pattern@token\_pattern}

\index{1fit\_transform@fit\_transform}

\index{1get\_feature\_names\_out@get\_feature\_names\_out}

\index{1vocabulary\_@vocabulary\_}

\index{1max\_features@max\_features}

\index{1min\_df@min\_df}

\index{1max\_df@max\_df}

\index{1OOV@OOV}

\index{1out-of-vocabulary@out-of-vocabulary}

\index{1CSR matrix@CSR matrix}

\index{1dense matrix@dense matrix}

\index{1load\_dataset@load\_dataset}

\index{1Yelp review full@Yelp review full}

\index{1pandas DataFrame@pandas DataFrame}

\index{0단어 가방@단어 가방}

\index{0엔그램@엔그램}

\index{0토큰 패턴@토큰 패턴}

\index{0어휘 사전@어휘 사전}

\index{0어휘 수@어휘 수}

\index{0어휘 밖 단어@어휘 밖 단어}

\index{0밀집 행렬@밀집 행렬}

\index{0데이터 샘플링@데이터 샘플링}


\section{변화추적표}

지금은 출발점이라 표가 한 줄뿐입니다. 장이 진행될수록 행이 누적됩니다.

\begin{booktable}{1장 변화추적표}{tab:ch01-tracking}
\begin{adjustbox}{max width=\textwidth}
\begin{tabular}{@{}llllll@{}}
\toprule
Ch & 모델 & 토크나이저 & Output Head & Activation & Loss \\
\midrule

\textbf{1 ← 여기} & (TF-IDF) & \inlinecode{TfidfVectorizer()} & --- & --- & --- \\
\bottomrule
\end{tabular}
\end{adjustbox}
\end{booktable}

전체 18개 장의 표는 \href{https://github.com/yoon-gu/neuqes-101\#챕터별-변화추적표}{루트 README.md}를 참고하세요.



\section{시작점에서}

이 장은 출발점입니다. 모델도 없고, 학습도 없고, loss도 없습니다. \textbf{텍스트를 숫자로 바꾸는 변환} 만 다룹니다.

왜 여기서 시작할까요? 자연어 모델은 결국 숫자 벡터를 다루는 함수입니다. BERT든 GPT든 입력 단계에서는 “텍스트 → 숫자” 변환이 반드시 필요합니다. 가장 단순한 변환부터 손에 익혀두면, 이후 장에서 BERT의 토크나이저와 임베딩이 어떻게 다른지 비교할 기준이 생깁니다.



\section{토크나이저 노트}

이 장의 토크나이저는 \textbf{\inlinecode{CountVectorizer} / \inlinecode{TfidfVectorizer} 자체} 입니다. 따로 토크나이저 객체가 분리돼 있지 않고, 벡터라이저 한 클래스가 아래 세 가지를 동시에 합니다.

\begin{enumerate}
\def\labelenumi{\arabic{enumi}.}
\tightlist
\item
  \textbf{분리(tokenize)}: 기본은 공백·구두점 기준 단어 단위 분리. 정규식 패턴 \texttt{(?u)\textbackslash{}b\textbackslash{}w\textbackslash{}w+\textbackslash{}b}(영숫자 2자 이상)이 기본값.
\item
  \textbf{어휘 학습(vocabulary)}: 학습 데이터에 등장한 토큰을 모두 모아 정수 인덱스로 매핑.
\item
  \textbf{OOV 처리}: 학습 어휘에 없는 단어는 \textbf{무시}합니다 (BERT의 \texttt{{[}UNK{]}}처럼 별도 토큰으로 보존하지 않음).
\end{enumerate}

조금 뒤 같은 문장 \inlinecode{"I\ love\ using\ Hugging\ Face!"} 가 어떻게 토큰화되는지 직접 확인합니다.

\begin{previewBox}{미리보기}
\textbf{다음 장(2장)}: 같은 \inlinecode{TfidfVectorizer}를 그대로 사용. 토크나이저는 변하지 않습니다.
\end{previewBox}



\Needspace{8\baselineskip}
\begin{lstlisting}
!pip install -q datasets scikit-learn pandas matplotlib
\end{lstlisting}

\begin{codeRead}
\textbf{1행}의 \inlinecode{!pip install -q datasets...}에서는 Colab 실행에 필요한 패키지를 설치합니다.
\end{codeRead}



\Needspace{11\baselineskip}
\begin{lstlisting}
import numpy as np
import pandas as pd
import matplotlib.pyplot as plt
from datasets import load_dataset
from sklearn.feature_extraction.text import CountVectorizer, TfidfVectorizer

plt.rcParams["axes.unicode_minus"] = False
\end{lstlisting}

\begin{codeRead}
\inlinecode{numpy, pandas, matplotlib, datasets, scikit-learn} 같은 주요 패키지를 불러옵니다.
\end{codeRead}



\section{실습: Yelp 리뷰 데이터 살펴보기}

\inlinecode{yelp\_review\_full}은 Yelp 식당 리뷰 65만 건에 1-5점 별점이 달린 데이터셋입니다 (라벨은 0-4로 저장됨). 학습 흐름을 가볍게 유지하기 위해 \textbf{5,000건만 무작위 샘플링} 합니다.



\Needspace{8\baselineskip}
\begin{lstlisting}
dataset = load_dataset("yelp_review_full")
dataset
\end{lstlisting}

\begin{codeRead}
\textbf{1행}의 \inlinecode{dataset = ...}에서는 이후 분석에서 사용할 값을 준비합니다. \textbf{2행}의 \inlinecode{dataset}에서는 앞 단계에서 만든 값을 바탕으로 다음 계산을 수행합니다.
\end{codeRead}

\noindent\textbf{출력.}
\begin{lstlisting}[style=bookoutput]
DatasetDict({
    train: Dataset({
        features: ['label', 'text'],
        num_rows: 650000
    })
    test: Dataset({
        features: ['label', 'text'],
        num_rows: 50000
    })
})
\end{lstlisting}
\par\vspace{0.9em}



\Needspace{10\baselineskip}
\begin{lstlisting}
SAMPLE_SIZE = 5000
ds = dataset["train"].shuffle(seed=42).select(range(SAMPLE_SIZE))
df = ds.to_pandas()

display(pd.DataFrame([{"metric": "sample_count", "value": len(df)}]))
df.head(3)
\end{lstlisting}

\begin{codeRead}
\textbf{1행}의 \inlinecode{SAMPLE\_SIZE = ...}에서는 이후 분석에서 사용할 값을 준비합니다. \textbf{2행}의 \inlinecode{ds = ...}에서는 전체 데이터에서 실습에 사용할 샘플만 추립니다. \textbf{3행}의 \inlinecode{df = ...}에서는 데이터셋을 표 형태로 바꿔 이후 셀에서 다루기 쉽게 합니다. \textbf{6행}의 \inlinecode{df.head(...)}에서는 앞 단계에서 만든 값을 바탕으로 다음 계산을 수행합니다.
\end{codeRead}

\noindent\textbf{출력.}
\begin{lstlisting}[style=bookoutput]
metric        value
-  ------------  -----
0  sample_count  5000
   label  text
-  -----  ------------------------
0  4      I stalk this truck. I...
1  2      who really knows if t...
2  4      I LOVE Bloom Salon......
\end{lstlisting}
\par\vspace{0.9em}



\Needspace{11\baselineskip}
\begin{lstlisting}
counts = df["label"].value_counts().sort_index()
labels = [f"{i+1} star" for i in counts.index]
plt.bar(labels, counts.values)
plt.title("Star rating distribution (sampled 5,000)")
plt.ylabel("Reviews")
plt.show()
display(counts.rename_axis("label").reset_index(name="count"))
\end{lstlisting}

\begin{codeRead}
\textbf{1행}의 \inlinecode{counts = ...}에서는 라벨별 샘플 개수를 집계합니다. \textbf{2행}의 \inlinecode{labels = ...}에서는 이후 분석에서 사용할 값을 준비합니다.
\end{codeRead}

\noindent\textbf{출력.}
\begin{lstlisting}[style=bookoutput]
<Figure size 640x480 with 1 Axes>
   label  count
-  -----  -----
0  0      1017
1  1      1027
2  2      960
3  3      1021
4  4      975
\end{lstlisting}
\par\vspace{0.9em}



\Needspace{8\baselineskip}
\begin{lstlisting}
df["len_words"] = df["text"].str.split().str.len()
df[["len_words"]].describe()
\end{lstlisting}

\begin{codeRead}
\textbf{1행}의 \inlinecode{df["len\_words"] = df["tex...}에서는 이후 분석에서 사용할 값을 준비합니다. \textbf{2행}의 \inlinecode{df[["len\_words"]].describe()}에서는 앞 단계에서 만든 값을 바탕으로 다음 계산을 수행합니다.
\end{codeRead}

\noindent\textbf{출력.}
\begin{lstlisting}[style=bookoutput]
len_words
count  5000.000000
mean    133.811400
std     119.787704
min       1.000000
25%      53.000000
50%     100.000000
75%     177.000000
max     977.000000
\end{lstlisting}
\par\vspace{0.9em}



\section{해부: CountVectorizer --- 텍스트를 “단어 횟수”로}

\inlinecode{CountVectorizer}는 가장 단순한 변환입니다.

\begin{quote}
“이 문서에 단어 X가 몇 번 나왔는가?”
\end{quote}

각 문서가 길이 V짜리 벡터로 변환됩니다 (V는 어휘 수). 대부분의 칸은 0이라 \textbf{희소(sparse)} 행렬로 저장합니다.



\Needspace{16\baselineskip}
\begin{lstlisting}
cv = CountVectorizer(max_features=10000)
X_count = cv.fit_transform(df["text"])
sparsity = 1 - X_count.nnz / (X_count.shape[0] * X_count.shape[1])

count_summary = pd.DataFrame([
    {"metric": "n_documents", "value": X_count.shape[0]},
    {"metric": "vocab_size", "value": X_count.shape[1]},
    {"metric": "non_zero_entries", "value": X_count.nnz},
    {"metric": "total_entries", "value": X_count.shape[0] * X_count.shape[1]},
    {"metric": "sparsity", "value": f"{sparsity:.2%}"},
])
display(count_summary)
\end{lstlisting}

\begin{codeRead}
\textbf{1행}의 \inlinecode{cv = ...}에서는 단어 횟수 기반 벡터라이저를 만듭니다. \textbf{2행}의 \inlinecode{X\_count = ...}에서는 텍스트나 라벨을 모델이 처리할 수 있는 수치 표현으로 변환합니다. \textbf{3행}의 \inlinecode{sparsity = ...}에서는 행렬에서 0인 원소의 비율을 계산합니다. \textbf{5--11행}의 \inlinecode{count\_summary = ...}에서는 결과를 표로 보기 좋게 정리합니다.
\end{codeRead}

\noindent\textbf{출력.}
\begin{lstlisting}[style=bookoutput]
metric            value
-  ----------------  --------
0  n_documents       5000
1  vocab_size        10000
2  non_zero_entries  405789
3  total_entries     50000000
4  sparsity          99.19%
\end{lstlisting}
\par\vspace{0.9em}



\Needspace{12\baselineskip}
\begin{lstlisting}
sample = "I love using Hugging Face!"
analyzer = cv.build_analyzer()

tokenization_preview = pd.DataFrame([
    {"item": "input", "value": repr(sample)},
    {"item": "tokens", "value": analyzer(sample)},
])
display(tokenization_preview)
\end{lstlisting}

\begin{codeRead}
\textbf{1행}의 \inlinecode{sample = ...}에서는 토큰화 예제로 사용할 문장을 정합니다. \textbf{2행}의 \inlinecode{analyzer = ...}에서는 벡터라이저 내부의 토큰화 규칙을 직접 호출할 함수로 꺼냅니다. \textbf{4--7행}의 \inlinecode{tokenization\_preview = ...}에서는 결과를 표로 보기 좋게 정리합니다.
\end{codeRead}

\noindent\textbf{출력.}
\begin{lstlisting}[style=bookoutput]
item    value
-  ------  ------------------------
0  input   'I love using Hugging...
1  tokens  [love, using, hugging...
\end{lstlisting}
\par\vspace{0.9em}



\textbf{관찰 포인트}

\begin{itemize}
\tightlist
\item
  모두 \textbf{소문자} 로 변환됩니다 (기본 \inlinecode{lowercase=True}).
\item
  구두점 \inlinecode{!}은 사라집니다 (정규식 패턴이 영숫자만 매칭).
\item
  \inlinecode{"I"} 같은 \textbf{단일 문자도 사라집니다} (기본 \inlinecode{token\_pattern}은 2자 이상만 인식).
\item
  학습 어휘에 없는 단어는 OOV로 \textbf{무시}됩니다 --- BERT처럼 \texttt{{[}UNK{]}}로 보존하지 않습니다.
\end{itemize}



\Needspace{18\baselineskip}
\begin{lstlisting}
vocab = cv.get_feature_names_out()
word_counts = np.asarray(X_count.sum(axis=0)).flatten()
top = np.argsort(word_counts)[::-1][:10]

vocab_summary = pd.DataFrame([{"metric": "vocab_size", "value": len(vocab)}])
vocab_preview = pd.DataFrame({"token": vocab[:20]})
top_count_words = pd.DataFrame({
    "token": vocab[top],
    "count": word_counts[top].astype(int),
})

display(vocab_summary)
display(vocab_preview)
display(top_count_words)
\end{lstlisting}

\begin{codeRead}
\textbf{1행}의 \inlinecode{vocab = ...}에서는 벡터라이저가 학습한 어휘 목록을 가져옵니다. \textbf{2행}의 \inlinecode{word\_counts = ...}에서는 열 방향으로 값을 더해 특성별 합계를 계산합니다. \textbf{3행}의 \inlinecode{top = ...}에서는 값이 큰 항목부터 볼 수 있도록 인덱스를 정렬합니다. \textbf{5행}의 \inlinecode{vocab\_summary = ...}에서는 결과를 표로 보기 좋게 정리합니다. \textbf{6행}의 \inlinecode{vocab\_preview = ...}에서는 결과를 표로 보기 좋게 정리합니다. \textbf{7--10행}의 \inlinecode{top\_count\_words = ...}에서는 결과를 표로 보기 좋게 정리합니다.
\end{codeRead}

\noindent\textbf{출력.}
\begin{lstlisting}[style=bookoutput]
metric      value
-  ----------  -----
0  vocab_size  10000
    token
0      00
1     000
2    00am
3    00pm
4      05
5      08
6      09
7      10
8     100
9    1000
10  100th
11    101
...
\end{lstlisting}
\par\vspace{0.9em}



\section{변형: TfidfVectorizer --- 흔한 단어의 영향력 줄이기}

\inlinecode{CountVectorizer}의 한계: \inlinecode{"the"}, \inlinecode{"and"} 같은 단어가 모든 리뷰에 많이 등장하니 횟수만으로는 문서 사이 차이를 잘 드러내지 못합니다.

\textbf{TF-IDF}는 두 항을 곱해 이 문제를 다룹니다.

\begin{equation}
\label{eq:ch01-01}
\text{tfidf}(t, d) = \underbrace{\text{tf}(t, d)}_{\text{문서 } d \text{에서 } t \text{의 빈도}} \cdot \underbrace{\log\frac{1 + N}{1 + \text{df}(t)}}_{\text{희귀도 가중치 (IDF)}}
\end{equation}

식~\eqref{eq:ch01-01}은 단어 빈도와 희귀도 가중치를 결합한 TF-IDF의 정의입니다.

\begin{itemize}
\tightlist
\item
  \inlinecode{tf}: 한 문서에서 단어가 얼마나 자주 나왔는가
\item
  \inlinecode{idf}: 그 단어가 \textbf{얼마나 적은 문서에 등장했는가} (모든 문서에 흔할수록 0에 가까워짐)
\end{itemize}

직관: “이 단어가 이 문서에서 자주 나오면서 동시에 다른 문서엔 흔하지 않다면, 이 문서를 특징짓는 단어”라는 가중치입니다.



\Needspace{8\baselineskip}
\begin{lstlisting}
tfidf = TfidfVectorizer(max_features=10000)
X_tfidf = tfidf.fit_transform(df["text"])
display(pd.DataFrame([{"object": "X_tfidf", "shape": X_tfidf.shape}]))
\end{lstlisting}

\noindent\textbf{출력.}
\begin{lstlisting}[style=bookoutput]
object   shape
-  -------  -------------
0  X_tfidf  (5000, 10000)
\end{lstlisting}
\par\vspace{0.9em}



\Needspace{22\baselineskip}
\begin{lstlisting}
doc_id = 0
review = df["text"].iloc[doc_id]

vocab_tf = tfidf.get_feature_names_out()
cv_row = np.asarray(X_count[doc_id].todense()).flatten()
tfidf_row = np.asarray(X_tfidf[doc_id].todense()).flatten()

top = np.argsort(tfidf_row)[::-1][:10]

review_preview = pd.DataFrame([{"review_preview_200_chars": review[:200] + "..."}])
top_tfidf_words = pd.DataFrame({
    "token": vocab_tf[top],
    "count": cv_row[top].astype(int),
    "tfidf": np.round(tfidf_row[top], 4),
})

display(review_preview)
display(top_tfidf_words)
\end{lstlisting}

\noindent\textbf{출력.}
\begin{lstlisting}[style=bookoutput]
review_preview_200_chars
0  I stalk this truck.  I've been to industrial p...
   token       count  tfidf
-  ----------  -----  ------
0  stalk       1      0.2418
1  pretend     1      0.2252
2  industrial  1      0.2212
3  farmer      1      0.2212
4  parks       1      0.2212
5  divine      1      0.2092
6  tech        1      0.2068
7  pride       1      0.2046
8  bowls       1      0.1988
9  gotta       1      0.1898
\end{lstlisting}
\par\vspace{0.9em}



\textbf{관찰}: 단순 횟수 기준 top 10에는 \inlinecode{the}, \inlinecode{and} 같은 흔한 단어가 위로 올라옵니다. TF-IDF 정렬에서는 그 문서를 특징짓는 명사·형용사가 상위로 올라오는 경향을 볼 수 있습니다.

\begin{quote}
\textbf{후속 논점}: 두 방식 모두 단어를 \emph{서로 독립} 으로 취급합니다. \inlinecode{"not\ bad"}(좋다는 뜻)와 \inlinecode{"bad"}(나쁘다는 뜻)을 구분하지 못합니다. BERT가 등장하는 Phase 1에서 이 한계가 어떻게 깨지는지 확인하게 됩니다.
\end{quote}



\section{이 장에 등장한 라이브러리}

\begin{booktable}{1장 새로 등장한 라이브러리}{tab:ch01-libraries}
\begin{adjustbox}{max width=\textwidth}
\begin{tabular}{@{}lll@{}}
\toprule
이름 & 한 줄 설명 & 다음 장에서 \\
\midrule

\inlinecode{datasets} & Hugging Face의 데이터셋 로딩 라이브러리 (Apache Arrow 기반) & 7장에서 깊게 본다 \\
\inlinecode{sklearn.feature\_extraction.text.CountVectorizer} & 횟수 벡터화 & 이후 장의 비교 기준 \\
\inlinecode{sklearn.feature\_extraction.text.TfidfVectorizer} & TF-IDF 벡터화 & 2-5장에서 입력으로 계속 사용 \\
\bottomrule
\end{tabular}
\end{adjustbox}
\end{booktable}



\section{체크포인트 질문}

\begin{enumerate}
\def\labelenumi{\arabic{enumi}.}
\tightlist
\item
  \inlinecode{CountVectorizer.fit\_transform(...)}이 만든 행렬의 shape가 \inlinecode{(N,\ V)}일 때, N과 V는 각각 무엇을 의미하나요?
\item
  sparsity가 99\%를 넘는 이유는 무엇인가요?
\item
  TF-IDF가 단순 횟수보다 “문서를 특징짓는 단어”를 더 잘 뽑아내는 이유는 IDF의 어느 부분에서 오나요?
\item
  \inlinecode{CountVectorizer}가 학습 어휘에 없는 단어를 만났을 때 어떻게 처리하나요? BERT의 처리 방식과는 무엇이 다른가요?
\end{enumerate}



\section{FAQ}

\begin{faqBox}{Q1. (실무) Yelp 데이터가 너무 커서 메모리에 안 올라가는데 어떻게 하나요?}

\inlinecode{datasets}는 Apache Arrow 메모리맵을 사용해서 65만 건 전체를 RAM에 올리지 않습니다. 인덱싱 시점에만 디스크에서 읽어 옵니다. 그래서 위 셀에서 \texttt{dataset{[}“train”{]}}을 가져와도 메모리는 거의 늘어나지 않습니다.

그래도 다운스트림(예: pandas 변환, sklearn fit) 단계에서 메모리가 부담되면 두 가지를 씁니다.

\Needspace{14\baselineskip}
\begin{lstlisting}
ds = dataset["train"].shuffle(seed=42).select(range(5000))

stream = load_dataset(
    "yelp_review_full",
    split="train",
    streaming=True,
)
for i, ex in enumerate(stream):
    if i >= 5000: break
    ...
\end{lstlisting}

\end{faqBox}

\begin{faqBox}{Q2. (실무) \inlinecode{CountVectorizer}와 \inlinecode{TfidfVectorizer} 중 뭘 써야 하나요?}

분류·검색 같은 일반 NLP 작업은 \textbf{TF-IDF가 기본 선택}입니다. 흔한 단어의 영향력을 자동으로 깎아주기 때문에 \inlinecode{LogisticRegression}, \inlinecode{LinearSVC} 같은 선형 모델과 잘 맞습니다.

\inlinecode{CountVectorizer}는 두 경우에 유리합니다.

\begin{itemize}
\tightlist
\item
  \inlinecode{MultinomialNB} (Naive Bayes): 빈도(count)를 가정한 모델이라 TF-IDF보다 Count가 더 잘 맞습니다.
\item
  “단어가 등장한 횟수 자체”가 의미 있는 분석(예: 단순 키워드 카운팅, 토픽 모델링 입력).
\end{itemize}

\end{faqBox}

\begin{faqBox}{Q3. (이론) TF-IDF의 IDF는 정확히 무슨 역할을 하나요?}

IDF(Inverse Document Frequency)는 \textbf{“이 단어가 얼마나 희귀한가”} 를 점수로 매깁니다. 식은 (sklearn 기본 옵션 기준):

\begin{equation}
\label{eq:ch01-02}
\text{idf}(t) = \log\frac{1 + N}{1 + \text{df}(t)} + 1
\end{equation}

식~\eqref{eq:ch01-02}은 문서 빈도로부터 IDF 값을 계산하는 방식입니다.

\begin{itemize}
\tightlist
\item
  \(N\): 전체 문서 수, \(\text{df}(t)\): 단어 \(t\)가 등장한 문서 수.
\item
  모든 문서에 등장하는 단어(\(\text{df}=N\)) → \(\log 1 = 0\), 거기에 \inlinecode{+1}이 더해져 IDF는 \textbf{1}.
\item
  한 문서에만 등장 → \(\log\frac{1+N}{2}\) 큰 값.
\end{itemize}

즉 흔한 단어는 IDF가 작아 TF-IDF 전체 값이 줄고, 드문 단어는 IDF가 커서 TF가 같아도 점수가 커집니다. “이 문서를 특징짓는 단어”를 골라내는 가중치가 IDF에서 옵니다.

\end{faqBox}

\begin{faqBox}{Q4. (이론) 어휘 수는 무엇을 기준으로 정해야 하나요? \inlinecode{max\_features=10000}은 어떻게 정한 값인가요?}

세 가지 트레이드오프가 있습니다.

\begin{enumerate}
\def\labelenumi{\arabic{enumi}.}
\tightlist
\item
  \textbf{너무 작으면}: 정보 손실. 핵심 단어가 어휘에서 잘려나감.
\item
  \textbf{너무 크면}: sparsity↑, 노이즈(오타·해시태그·고유명사 1회 등장 단어)도 같이 학습.
\item
  \textbf{모델 학습 시간/메모리}: V가 커지면 선형 모델 가중치도 V만큼 커짐.
\end{enumerate}

실무에선 보통 두 가지를 조합합니다.

\Needspace{9\baselineskip}
\begin{lstlisting}
TfidfVectorizer(
    max_features=10000,
    min_df=5,
    max_df=0.9,
)
\end{lstlisting}

5,000건짜리 영어 문서 데이터에선 5K-30K가 흔한 출발점입니다. \inlinecode{10000}은 학습 흐름을 빠르게 가져가기 위한 보수적 설정이고, 실험으로 조정하면 됩니다.

\end{faqBox}

\begin{faqBox}{Q5. (실무) sklearn에서 한국어도 처리되나요? 다음 장에서도 그대로 동작할까요?}

\textbf{기본 정규식 토크나이저로도 동작은 합니다} (한국어도 영숫자 패턴으로 잡힘). 다만 한국어는 조사 때문에 같은 어근이 다른 토큰으로 쪼개져 어휘가 폭증합니다 --- 예: \inlinecode{학교}, \inlinecode{학교는}, \inlinecode{학교가}, \inlinecode{학교를}이 전부 다른 토큰.

실무에선 형태소 분석기를 토크나이저로 끼워 넣습니다.

\Needspace{8\baselineskip}
\begin{lstlisting}
from konlpy.tag import Mecab
tokenizer = Mecab().morphs

TfidfVectorizer(tokenizer=tokenizer, token_pattern=None)
\end{lstlisting}

이 커리큘럼에서는 Phase 2(13-16장, 한국어)에서 \inlinecode{klue/bert-base}의 한국어 WordPiece 토크나이저를, Phase 3(18장)에서 형태소 기반 워드레벨 토크나이저를 직접 다룹니다.

\end{faqBox}

\begin{faqBox}{Q6. (이론) sparse 행렬이 dense 행렬보다 메모리에 유리한 이유는 무엇인가요? \inlinecode{.toarray()}로 바꾸면 왜 메모리가 폭발할 수 있나요?}

shape \inlinecode{(5000,\ 10000)} 행렬을 dense(\inlinecode{float64})로 만들면 \inlinecode{5000\ ×\ 10000\ ×\ 8\ byte\ ≈\ 400MB}입니다. 칸 대부분이 0인데 그 0까지 다 저장합니다.

sparse(CSR) 행렬은 0이 아닌 칸만 \inlinecode{(값,\ 열\ 인덱스)}로 저장합니다. nnz가 50만이면 대략 \inlinecode{500000\ ×\ (8\ +\ 4)\ ≈\ 6MB} --- 70배 가까이 절약됩니다.

\Needspace{9\baselineskip}
\begin{lstlisting}
memory_check = pd.DataFrame([
    {"layout": "sparse", "nbytes_mb": (X_count.data.nbytes + X_count.indices.nbytes + X_count.indptr.nbytes) / 1e6},
    {"layout": "dense", "nbytes_mb": (X_count.shape[0] * X_count.shape[1] * 8) / 1e6},
])
display(memory_check.round(1))
\end{lstlisting}

\inlinecode{X\_count.toarray()}를 호출하면 dense로 풀어 그 400MB짜리 배열을 메모리에 올립니다. Yelp 5,000건 정도는 버티지만, 50,000건 + max\_features 50,000으로 키우면 \inlinecode{50000\ ×\ 50000\ ×\ 8\ =\ 20GB}가 되어 Colab이 그 자리에서 중단될 수 있습니다. \textbf{sklearn 모델 대부분(LinearRegression, LogisticRegression, LinearSVC 등)은 sparse 입력을 그대로 받기 때문에 굳이 \inlinecode{.toarray()}로 바꿀 일이 없습니다.}
\end{faqBox}



\section{오류 실험 (선택)}

다음 코드를 돌려보고 결과를 비교해보세요. 어디가 달라졌나요?

\Needspace{11\baselineskip}
\begin{lstlisting}
cv2 = CountVectorizer(
    max_features=10000,
    lowercase=False,
    token_pattern=r"\b\w+\b",
)
X2 = cv2.fit_transform(df["text"])
display(pd.DataFrame([{"metric": "vocab_size", "value": len(cv2.get_feature_names_out())}]))
\end{lstlisting}

힌트: \inlinecode{lowercase=False}로 두면 \inlinecode{"good"}과 \inlinecode{"Good"}이 같은 토큰일까요, 다른 토큰일까요? 어휘 수가 변하는 방향을 예측해보세요.



\begin{previewBox}{미리보기: 다음 장}

\textbf{2장. sklearn Regression --- 시작점}

\begin{itemize}
\tightlist
\item
  \inlinecode{LinearRegression}으로 별점(1-5)을 회귀합니다.
\item
  활성화 함수도 없이 출력값을 그대로 사용 --- 음수도, 5보다 큰 값도 나올 수 있습니다.
\item
  다음 장의 첫 Loss 등장: \inlinecode{MSELoss} (sklearn: squared error).
\end{itemize}
\end{previewBox}
